\chapter{Lagrangian Relaxation e Column Generation}
\label{ch:metodo_lagrangian}

\epigraph{``If you cannot solve a problem, relax it.''}{adagio
della ricerca operativa, attribuito a vari autori}

\lettrine[lines=3]{C}{i sono} problemi che, pur essendo
risolvibili in linea di principio da CP-SAT, sono \emph{cos\`i
grandi} che il solver impiega ore per produrre una soluzione
mediocre. Una scuola con 200 classi, 400 docenti, 12 indirizzi
non sta nei tempi di Phase~B di un ciclo standard. Per questi
casi $\pi$Tantum offre due tecniche di \emph{rilassamento e
decomposizione} di derivazione classica: il rilassamento
lagrangiano e la generazione di colonne. Entrambe sono di
default \emph{disattivate}, da abilitare esplicitamente da
chi gestisce la pipeline.

\section{Rilassamento lagrangiano: dualizzare i ponti}

\begin{storiabox}
Il rilassamento lagrangiano nasce nella ricerca operativa
degli anni '60--'70. Held e Karp (1971)~\parencite{heldkarp1971}
lo applicarono al \emph{Traveling Salesman Problem} ottenendo
risultati eccellenti su istanze grandi. Geoffrion
(1974)~\parencite{geoffrion1974} formaliz\`o l'approccio per
la programmazione intera generale; Fisher
(1981)~\parencite{fisher1981} ne fece il manuale operativo
per chi voleva applicarlo nella pratica industriale.
\end{storiabox}

\begin{ideabox}
\textbf{L'idea con un'analogia}. Hai un problema enorme che
non riesci a risolvere. Decidi di \emph{dimenticare} alcuni
vincoli scomodi (che ``legano'' parti diverse del problema)
e di \emph{punirli} con una multa proporzionale alla loro
violazione. Risolvi il problema rilassato. Misuri quanto i
vincoli scomodi sono violati. Aggiusti la multa: dove sono
violati di pi\`u, alzi la multa; dove sono soddisfatti con
margine, la abbassi. Iteri.

Si dimostra che, sotto ipotesi tecniche, le multe convergono
a un valore tale che la soluzione del problema rilassato
soddisfa anche i vincoli ``dimenticati'' --- ottenendo
quindi una soluzione del problema originale.
\end{ideabox}

\subsection{Formalizzazione}

Supponi un problema di minimo:
\[
  \min_{x \in X} f(x) \quad \text{soggetto a } g(x) \leq 0
\]
dove $X$ \`e un insieme di soluzioni ``facili'' (es. unione
di sotto-problemi indipendenti) e $g(x) \leq 0$ rappresenta
i vincoli ``difficili'' (i ponti fra sotto-problemi).

\begin{defbox}
La \emph{Lagrangiana}\index{Lagrangiana} \`e:
\[
  L(x, \lambda) = f(x) + \lambda^\top g(x)
\]
con $\lambda \geq 0$ vettore dei \emph{moltiplicatori di
Lagrange}.

Il \emph{problema rilassato} \`e:
\[
  q(\lambda) = \min_{x \in X} L(x, \lambda)
\]
La funzione $q(\lambda)$ \`e detta \emph{funzione duale}.
\end{defbox}

\textbf{Come si legge}: per ogni valore del vettore $\lambda$,
risolvi un problema senza i vincoli scomodi $g(x) \leq 0$
ma con una penalit\`a $\lambda^\top g(x)$ aggiunta
all'obiettivo. Se $g(x) > 0$ (vincolo violato), la penalit\`a
sale; se $g(x) < 0$ (vincolo soddisfatto con margine), la
penalit\`a scende (il termine \`e negativo). Questo ``incentiva''
$x$ a soddisfare i vincoli senza forzarli.

\subsection{Subgradient ascent: aggiornare i moltiplicatori}

Si dimostra che $q(\lambda)$ \`e concava in $\lambda$ e che
$\max_{\lambda \geq 0} q(\lambda) \leq \min_{x: g(x) \leq 0} f(x)$
(weak duality). L'algoritmo classico aggiorna $\lambda$ con
un passo di \emph{subgradient ascent}\index{subgradient}:
\[
  \lambda_{k+1} = \max\left(0,\ \lambda_k + \alpha_k g(x_k)\right)
\]
dove $\alpha_k > 0$ \`e il \emph{passo}, tipicamente
$\alpha_k = \alpha_0 / (1 + k)$ in $\pi$Tantum.

\textbf{Come si legge}: dopo aver risolto il rilassato e
trovato $x_k$, calcolo quanto i vincoli sono violati ($g(x_k)$).
Aggiorno $\lambda$ aggiungendoci un multiplo di $g(x_k)$:
i moltiplicatori associati a vincoli pi\`u violati salgono di
pi\`u. Tronco a zero per mantenere $\lambda \geq 0$.

\subsection{Cosa rilassare nel timetabling?}

In $\pi$Tantum la decomposizione lagrangiana si applica
\emph{dopo} la decomposizione spettrale: si rilassano i
\emph{ponti inter-cluster}\index{ponte inter-cluster}, cio\`e
i vincoli che legano lezioni di cluster diversi (tipicamente
docenti che insegnano a cavallo fra due cluster, oppure
vincoli logici DSL che attraversano cluster).

Ogni cluster diventa un sotto-problema CP-SAT indipendente
con un costo aggiuntivo $\lambda_b \cdot (\text{violazione del
ponte } b)$. Iterando il subgradient, il sistema converge a
un assegnamento dei ponti coerente con tutti i cluster.

\subsection{Esempio in miniatura}

\begin{esempiobox}
Due cluster: A (3 classi) e B (3 classi). Un docente, Conti,
deve insegnare 2 ore in cluster A (storia in 1A) e 2 ore in
cluster B (storia in 4B). Vincolo ponte: ``Conti non in due
slot contemporanei''.

Senza rilassamento: Phase~B risolve A e B insieme, vincolo
HARD globale, propagazione attraversa i cluster (lento).

Con rilassamento: $\lambda$ inizia a 0. Risolvo A da solo
(mette Conti in lun-1, lun-2). Risolvo B da solo (mette
Conti in lun-1, lun-2). Conflitto: Conti in lun-1 e lun-2 in
entrambi i cluster.

Aggiorno $\lambda \leftarrow \alpha \cdot 2 = 0.2$
(violazione = 2 slot duplicati).

Iter 2: A ora penalizza Conti in lun-1, lun-2 con costo 0.2
ciascuno; trova Conti in mar-1, mar-2 (costo SOFT pi\`u alto
ma niente penalit\`a). B mantiene Conti in lun-1, lun-2.
Conflitto risolto, ma A ha pagato un piccolo costo SOFT.

Dopo 5--10 iterazioni il sistema trova un equilibrio in cui
Conti \`e in slot diversi nei due cluster e nessun vincolo
globale \`e violato.
\end{esempiobox}

\subsection{Perch\'e non sempre?}

\begin{avvertenzabox}
Lagrangian Relaxation \`e potente ma costoso da
implementare bene:
\begin{itemize}
\item Richiede il \emph{tuning} del passo $\alpha_0$.
  Troppo grande $\to$ oscillazioni; troppo piccolo $\to$
  convergenza lentissima.
\item Pu\`o \emph{non convergere} se il problema duale ha
  un \emph{duality gap} grosso (frequente nei problemi
  combinatori interi).
\item La soluzione finale del rilassato \emph{non} \`e
  necessariamente fattibile per il problema originale:
  serve un passo di ``ricostruzione'' (in $\pi$Tantum: un
  giro di SA per chiudere i ponti residui).
\end{itemize}

Per queste ragioni $\pi$Tantum lo offre come \emph{opzione
avanzata}, non come default.
\end{avvertenzabox}

\section{Column Generation: il catalogo immobiliare}

\begin{storiabox}
La generazione di colonne (\emph{Column Generation}) nasce
nel 1960 con Dantzig e Wolfe~\parencite{dantzig1960}.
L'idea: in molti grandi problemi di programmazione lineare,
\emph{la maggior parte delle variabili non serve}; sono
attive solo poche. Invece di costruire il problema con
tutte le variabili, lo si costruisce con poche e si
\emph{generano nuove colonne} (= nuove variabili) solo
quando servono.
\end{storiabox}

\begin{ideabox}
\textbf{Analogia}. Sei un'agenzia immobiliare con un milione
di appartamenti in catalogo. Un cliente cerca un trilocale
in zona centro. Non gli passi tutto il catalogo: gli
proponi 3--4 appartamenti che pensi gli vadano bene.
Se non gli piacciono, ne tiri fuori altri 3--4 con
caratteristiche diverse. Itera finch\'e il cliente \`e
soddisfatto. Hai usato il catalogo intero ma esaminandone
solo una frazione.

Column generation fa lo stesso: ha un \emph{master problem}
(piccolo, con poche variabili) e uno o pi\`u \emph{sub-problem}
che \emph{generano} nuove variabili da aggiungere al master,
basandosi sui prezzi duali correnti del master.
\end{ideabox}

\subsection{Schema operativo}

\begin{enumerate}
\item Costruisci il master problem con un sotto-insieme
  iniziale di variabili (colonne).
\item Risolvi il master come LP. Ottieni i prezzi duali
  $\pi$ delle restrizioni.
\item Per ogni sotto-problema, cerca una colonna con
  \emph{costo ridotto negativo} ($c_j - \pi^\top a_j < 0$).
  Se ne trovi, aggiungila al master.
\item Se nessun sotto-problema produce una colonna
  vantaggiosa, fermati: il master corrente \`e ottimo per
  l'LP rilassato.
\item (Per problemi interi) lancia un branch-and-bound
  sopra (\emph{branch-and-price}).
\end{enumerate}

\subsection{Cosa generano le colonne nel timetabling?}

In $\pi$Tantum la formulazione \`e: il master \`e un
\emph{set covering} dove ogni colonna rappresenta un
\emph{pattern settimanale completo} di un docente (la
sequenza di tutti i suoi slot in tutte le sue classi). Il
master sceglie un pattern per ogni docente in modo da
coprire tutte le ore-cattedra delle classi. I sub-problem
generano nuovi pattern docenti, uno per docente, sulla base
dei prezzi duali.

Il vantaggio: il numero di pattern possibili per docente \`e
combinatorialmente esplosivo (centinaia di milioni), ma il
master ne usa solo qualche decina. Il sotto-problema \`e
relativamente piccolo (un docente alla volta) e CP-SAT lo
risolve in pochi millisecondi.

\subsection{Granularit\`a configurabile del sotto-problema}

Lo schema di Dantzig-Wolfe non impone che il pattern sia
\emph{per docente}. La granularit\`a -- ovvero l'unit\`a di
decomposizione su cui il sotto-problema CP-SAT genera nuovi
pattern -- \`e una scelta di modellazione che $\pi$Tantum
espone come opzione. Quattro granularit\`a sono previste,
ognuna adatta a una taglia o a una struttura diversa della
scuola.

La granularit\`a \emph{per docente} (\texttt{teacher}) \`e quella
gi\`a descritta sopra: un pattern \`e una proposta di
settimana per un singolo docente, e si genera un sub per
ogni docente. Funziona meglio nelle scuole piccole, dove il
numero di docenti \`e dominante e i loro cataloghi
settimanali coprono bene lo spazio.

La granularit\`a \emph{per classe} (\texttt{class}) inverte
la prospettiva: un pattern \`e la settimana di lezioni di una
classe, con i suoi docenti gi\`a assegnati. Si genera un sub
per ogni classe. Scala meglio della prima quando il numero
di classi \`e moderato (30-80) ma i cataloghi per docente
esplodono.

La granularit\`a \emph{per giorno} (\texttt{day}) \`e una
proiezione orizzontale: un pattern \`e una giornata intera
della scuola, con tutti i suoi vincoli intra-day. Il sub
\`e quindi un per giorno della settimana. Raramente \`e la
scelta migliore in BnP puro -- la decomposizione temporale
del cap.~\ref{ch:metodo_spettrale} la sfrutta in modo pi\`u
diretto -- ma resta selezionabile per esperimenti.

La granularit\`a \emph{per indirizzo} (\texttt{curriculum})
\`e quella naturalmente adatta alle scuole italiane di
medie-grandi dimensioni. Un pattern \`e la settimana di
tutte le classi di un certo indirizzo (Scientifico,
Classico, Linguistico, ITIS Informatica, eccetera), con i
docenti che principalmente vi insegnano. Il sub \`e quindi
un per ciascun indirizzo. I docenti-ponte (chi insegna in
pi\`u indirizzi: tipicamente inglese, scienze motorie,
religione) sono modellati come constraint del master che
coordinano tra patterns di indirizzi diversi attraverso i
duali. \`E lo schema classico di
Vanderbeck~\parencite{vanderbeck2010} per Dantzig-Wolfe
applicato a problemi con \emph{struttura bloccata}
(\emph{block-angular}): il problema globale si decompone in
blocchi quasi-indipendenti (gli indirizzi) collegati da
poche variabili globali (i ponti).

L'interfaccia espone una opzione \emph{Auto} che il backend
risolve in base al numero di classi della scuola attiva:
sotto le trenta classi suggerisce \texttt{teacher},
fra trenta e ottanta \texttt{class}, sopra le ottanta con
indirizzi ben definiti \texttt{curriculum}. Per il profilo
MEGA (cento classi) l'auto-detect propone proprio
\texttt{curriculum}, dato che gli otto indirizzi della
scuola hanno pool di docenti largamente disgiunti.

\subsection{Stato attuale e limiti pratici}

\begin{avvertenzabox}
La versione attuale di $\pi$Tantum esegue una variante
\emph{iterative-diversified} della column generation: master
LP iterativo (\texttt{scipy.linprog} con backend HiGHS),
arricchimento del catalogo con varianti casuali ad ogni
iterazione, recupero intero per docente con il pattern di
peso LP massimo, e \emph{completion solver} day-by-day che
riempie eventuali gap residui invocando direttamente
\texttt{cpsat\_v2\_timetable}. Il risultato \`e
HARD-feasibile in tutti i casi misurati su \texttt{small}.

Il vero \emph{branch-and-price} con sub-CP-SAT guidato dai
duali per ciascuna delle quattro granularit\`a, e con
branching di Ryan-Foster per la ricerca dell'integer feasible,
\`e progettato e documentato in
\texttt{docs/optimization\_strategies.md} ma non \`e ancora
implementato. La stima realistica \`e di cinque-dieci
giorni di engineering OR focalizzato. La struttura del
modulo \texttt{experiments/column\_generation.py} \`e
pronta per ospitare un secondo percorso
(\texttt{mode="branch-and-price"}) accanto all'attuale
\texttt{mode="iterative-diversified"}, ma il refactor del
constraint del master e' la parte invasiva.

Default: \texttt{enable\_cg=False} nella pipeline.
\end{avvertenzabox}

\section{Connessioni}

\begin{connessionebox}
\begin{itemize}
\item \textbf{Decomposizione spettrale}
  (cap.~\ref{ch:metodo_spettrale}): naturale alleato del
  rilassamento lagrangiano (i cluster diventano i
  sotto-problemi, i ponti inter-cluster diventano i vincoli
  da dualizzare).
\item \textbf{CP-SAT} (cap.~\ref{ch:metodo_cpsat}): risolve i
  sotto-problemi della decomposizione lagrangiana e i
  sub-problem della column generation.
\item \textbf{Metaeuristiche}
  (cap.~\ref{ch:metodo_metaeuristiche}): un giro di SA dopo
  Lagrangian per chiudere i ponti residui.
\end{itemize}
\end{connessionebox}

\begin{biblioparvabox}
\begin{itemize}
\item \cite{geoffrion1974}, \cite{fisher1981}: i due
  riferimenti classici sul rilassamento lagrangiano per IP.
\item \cite{heldkarp1971}: l'applicazione storica al TSP.
\item \cite{dantzig1960}: il paper originale sulla
  decomposizione di Dantzig-Wolfe e generazione di
  colonne.
\item \cite{desrosiers2005}: \emph{A Primer in Column
  Generation}, manuale moderno e accessibile.
\item \cite{vanderbeck2010}: trattamento sistematico delle
  decomposizioni \emph{block-angular} per programmazione
  intera, su cui poggia la granularit\`a per indirizzo.
\item \cite{ryanfoster1981}: lo schema di branching che
  porta il loro nome, riferimento standard nei moderni
  branch-and-price.
\item \cite{wolsey1998}: capp.~10--11 sui rilassamenti
  e le decomposizioni in IP.
\end{itemize}
\end{biblioparvabox}
